\documentclass[12pt]{ximera}
\title{Exam 3 (2022)}
\begin{document}
\begin{abstract}
An archived exam on Chapter 3 and the start of Chapter 4: lone-divider bids, dividing a pizza, sealed bids with positive and negative assets, the method of markers, and Hamilton and Jefferson apportionments.
\end{abstract}
\maketitle

\section*{Exam 3 (2022)}

This exam allowed one handwritten sheet of notes.

\begin{problem}
Four brothers A, B, C, D are dividing their real estate empire using the lone-divider
method. The divider made four shares $s_1,\dots,s_4$ that are equal in his view, and the
others submitted their bids:
\begin{center}
\begin{tabular}{r|cccc}
 & A & B & C & D\\\hline
Acceptable & $\set{s_3}$ & $\set{s_1,s_2}$ & $\set{s_1,s_2,s_3,s_4}$ & $\set{s_2,s_3}$
\end{tabular}
\end{center}

\begin{prompt}
\textbf{(a)} Which brother was the divider? $\answer{C}$
\begin{feedback}
The divider finds every share acceptable, and only C does.
\end{feedback}
\end{prompt}

\begin{prompt}
\textbf{(b)} Which share does each brother get? Enter share numbers.

A: $\answer{3}$ \quad B: $\answer{1}$ \quad C: $\answer{4}$ \quad D: $\answer{2}$
\begin{feedback}
A must get $s_3$. That leaves only $s_2$ for D, and then only $s_1$ for B. C, the
divider, gets $s_4$.
\end{feedback}
\end{prompt}
\end{problem}

\begin{problem}
Three friends A, B, C share a pizza that is one quarter each cheese, salami, veggie, and
supreme, using the lone-divider method with A as divider. A likes cheese and salami
equally, and veggie as much as cheese and salami together, but gives the supreme part a
value of \$0. The whole pizza is worth \$12 to A.

\begin{prompt}
\textbf{(a)} What is each part worth to A, and how much is each of A's shares worth?

Cheese $\$\answer{3}$ \quad Salami $\$\answer{3}$ \quad Veggie $\$\answer{6}$ \quad
Supreme $\$\answer{0}$ \qquad Each share: $\$\answer{4}$
\begin{feedback}
If cheese and salami are worth $x$ each, veggie is worth $2x$, so $4x=12$ and $x=3$.
Each of three equal shares is worth $\$12\div3=\$4$.
\end{feedback}
\end{prompt}

\begin{prompt}
\textbf{(b)} Describe one way for A to make three shares.
\begin{freeResponse}
\end{freeResponse}
\begin{feedback}
Answers vary. One way: start with all the cheese ($\$3$) and add $\$1$ of salami, which
is $\frac13$ of it: $s_1=C+\frac13S$. The remaining $\frac23S$ is worth $\$2$, so add
$\$2$ of veggie, which is $\frac13$ of it: $s_2=\frac23S+\frac13V$. The rest is
$s_3=\frac23V+Su$, automatically worth $\$4$.
\end{feedback}
\end{prompt}

\begin{prompt}
\textbf{(c)} B values salami and supreme equally and doesn't like cheese or veggie at
all. For the shares in the sample answer to (b), what is each share worth to B (as
dollars of a \$12 pizza)?

$s_1$: $\$\answer{2}$ \quad $s_2$: $\$\answer{4}$ \quad $s_3$: $\$\answer{6}$
\begin{feedback}
To B, salami and supreme are worth $\$6$ each. $s_1=\frac13\cdot6=2$,
$s_2=\frac23\cdot6=4$, and $s_3=6$. B finds $s_2$ and $s_3$ acceptable.
\end{feedback}
\end{prompt}
\end{problem}

\begin{problem}
Four siblings divide the assets of a toothbrush factory by sealed bids. Bids are in
thousands of dollars.
\begin{center}
\begin{tabular}{r|rrrr}
 & Adi & Beatrice & Chris & Deanna\\\hline
Office building & 500 & 350 & 400 & 450\\
Unsold product & 350 & 360 & 400 & 390\\
Warehouse & 150 & 130 & 200 & 240
\end{tabular}
\end{center}

\begin{prompt}
\textbf{(a)} Who gets each item?

Office: $\answer{Adi}$ \quad Product: $\answer{Chris}$ \quad Warehouse: $\answer{Deanna}$
\begin{feedback}
Each item goes to its highest bidder; Beatrice gets no item.
\end{feedback}
\end{prompt}

\begin{prompt}
\textbf{(b)} Find the initial settlement for each sibling (positive = pays in, negative =
receives).

Adi $\answer{250}$ \quad Beatrice $\answer{-210}$ \quad Chris $\answer{150}$ \quad
Deanna $\answer{-30}$
\begin{feedback}
\[
\begin{array}{r|rrrr}
 & \text{Adi} & \text{Beatrice} & \text{Chris} & \text{Deanna}\\\hline
\text{Total} & 1000 & 840 & 1000 & 1080\\
\text{Fair share} & 250 & 210 & 250 & 270\\
\text{Items} & 500 & 0 & 400 & 240\\
\text{Items}-\text{fair} & 250 & -210 & 150 & -30
\end{array}
\]
\end{feedback}
\end{prompt}

\begin{prompt}
\textbf{(c)} How much surplus is left? $\answer{160}$ thousand, or $\answer{40}$
thousand per sibling
\begin{feedback}
$250+150-210-30=160$, split four ways.
\end{feedback}
\end{prompt}
\end{problem}

\begin{problem}
Two housemates must clean before returning the keys, and they use sealed bids with
negative assets. Their estimates of each chore's cost, in dollars:
\begin{center}
\begin{tabular}{r|rr}
 & Anna & Billy\\\hline
Clean windows & 23 & 35\\
Vacuum carpet & 15 & 12\\
Clean bathroom & 30 & 28\\
Take out trash & 22 & 25
\end{tabular}
\end{center}

\begin{prompt}
\textbf{(a)} Which chores does Anna do?
\begin{selectAll}
\choice[correct]{Windows}
\choice{Carpet}
\choice{Bathroom}
\choice[correct]{Trash}
\end{selectAll}
\begin{feedback}
Each chore goes to whoever minds it least: Anna does windows and trash, Billy the carpet
and bathroom.
\end{feedback}
\end{prompt}

\begin{prompt}
\textbf{(b)} Initial settlement (positive = pays): Anna $\answer{0}$ \quad Billy
$\answer{10}$
\begin{feedback}
Anna's total is $90$, so her fair share of the work is $45$, exactly what her chores
cost her ($23+22$). Billy's total is $100$, a fair share $50$, but his chores cost only
$12+28=40$, so he pays the $\$10$ difference.
\end{feedback}
\end{prompt}

\begin{prompt}
\textbf{(c)} What is the surplus? $\$\answer{10}$
\begin{feedback}
The $\$10$ Billy pays in is split $\$5$ each: Anna ends up receiving $\$5$, and Billy
pays a net $\$5$.
\end{feedback}
\end{prompt}
\end{problem}

\begin{problem}
Three players divide $14$ pieces of candy, lined up in a row, by the method of markers.
A's fair shares are $(1\text{--}4)$, $(5\text{--}8)$, $(9\text{--}14)$; B's are
$(1\text{--}7)$, $(8\text{--}11)$, $(12\text{--}14)$; and C's are $(1\text{--}3)$,
$(4\text{--}8)$, $(9\text{--}14)$.

\begin{prompt}
\textbf{(a)} How is the candy allocated? Enter the first and last piece.

C gets $\answer{1}$--$\answer{3}$ \quad A gets $\answer{5}$--$\answer{8}$ \quad
B gets $\answer{12}$--$\answer{14}$
\begin{feedback}
The leftmost first marker is C's (after piece $3$), so C takes $1$--$3$. Of the
remaining second markers, A's (after $8$) comes before B's (after $11$), so A takes
$5$--$8$. B gets its last share, $12$--$14$.
\end{feedback}
\end{prompt}

\begin{prompt}
\textbf{(b)} Which pieces are left over?
\begin{selectAll}
\choice[correct]{4}
\choice{8}
\choice[correct]{9}
\choice[correct]{10}
\choice[correct]{11}
\choice{12}
\end{selectAll}
\begin{feedback}
Pieces $4$ and $9$--$11$ form the surplus.
\end{feedback}
\end{prompt}
\end{problem}

\begin{problem}
A college reorganizes into four departments: Animal Science, Biology, Chemistry, and Math
(A, B, C, M). It will apportion $30$ administrative staffers by student enrollment using
Hamilton's method.
\begin{center}
\begin{tabular}{r|rrrrr}
 & A & B & C & M & Total\\\hline
Enrollment & 3700 & 1400 & 1600 & 2300 & 9000
\end{tabular}
\end{center}

\begin{prompt}
\textbf{(a)} Standard divisor: $\answer{300}$
\end{prompt}

\begin{prompt}
\textbf{(b)} Standard quotas: A $\answer[tolerance=0.005]{12.33}$, B
$\answer[tolerance=0.005]{4.67}$, C $\answer[tolerance=0.005]{5.33}$, M
$\answer[tolerance=0.005]{7.67}$
\end{prompt}

\begin{prompt}
\textbf{(c)} Hamilton apportionment: A $\answer{12}$, B $\answer{5}$, C $\answer{5}$,
M $\answer{8}$
\begin{feedback}
Lower quotas $12,4,5,7$ total $28$, so $2$ staffers remain. The largest residues are
B and M ($0.67$ each).
\end{feedback}
\end{prompt}

\begin{prompt}
\textbf{(d)} With $31$ staffers the standard quotas become $12.7$, $4.8$, $5.5$, and $7.9$.
Which department gets the extra staffer, compared with part (c)? $\answer{A}$
\begin{feedback}
Lower quotas still total $28$, now with $3$ staffers to hand out. The largest residues
are M ($0.9$), B ($0.8$), and A ($0.7$), giving A $13$, B $5$, C $5$, M $8$. Only A
changed.
\end{feedback}
\end{prompt}
\end{problem}

\begin{problem}
The college now uses Jefferson's method for the $30$ staffers instead. Modified quotas for
several divisors:
\begin{center}
\begin{tabular}{r|rrrr}
Divisor & A & B & C & M\\\hline
$290$ & 12.76 & 4.83 & 5.52 & 7.93\\
$287$ & 12.89 & 4.88 & 5.57 & 8.01\\
$285$ & 12.98 & 4.91 & 5.61 & 8.07\\
$281$ & 13.17 & 4.98 & 5.69 & 8.19\\
$280$ & 13.21 & 5.00 & 5.71 & 8.21
\end{tabular}
\end{center}

\begin{prompt}
\textbf{(a)} Which divisor works? $\answer{281}$
\begin{feedback}
Rounding down: $290\to28$, $287\to29$, $285\to29$, $281\to30$, $280\to31$.
\end{feedback}
\end{prompt}

\begin{prompt}
\textbf{(b)} Jefferson apportionment: A $\answer{13}$, B $\answer{4}$, C $\answer{5}$,
M $\answer{8}$
\begin{feedback}
Compared with Hamilton, Jefferson moves a staffer from the smaller department B to the
largest, A, which is typical of its bias toward large states.
\end{feedback}
\end{prompt}
\end{problem}

\end{document}
